\subsection{Cache và quản lý độ mới}
\label{subsec:db_cache}

Yoca sử dụng hai cách xác định độ mới. Với snapshot như metadata, market overview hoặc kết quả phân tích, TTL được tính từ thời điểm cập nhật; bản ghi còn hiệu lực được trả trực tiếp, còn bản ghi hết hạn được thay thế sau khi phản hồi provider đã qua kiểm tra schema. Với dữ liệu lịch sử, coverage metadata ghi nhận khoảng thời gian đã đồng bộ. Khi người dùng yêu cầu một kỳ rộng hơn, service chỉ lấy đoạn còn thiếu và hợp nhất với dữ liệu đã có.

Các bản ghi chuỗi thời gian được upsert theo address cùng timestamp hoặc signature. Cách này vừa chống trùng khi provider trả lại dữ liệu cũ, vừa cho phép cập nhật đoạn cuối của lịch sử mà không tải lại toàn bộ ví. Metadata, market snapshot, chart, pool và holder được tách thành các read model riêng vì nguồn và chu kỳ làm mới khác nhau; một phần hết hạn vì vậy không buộc hệ thống gọi lại mọi endpoint liên quan.

\subsection{Toàn vẹn và bảo mật dữ liệu}
\label{subsec:db_toan_ven}

Primary key tổng hợp được dùng cho dữ liệu chuỗi thời gian và lịch sử, chẳng hạn address kết hợp timestamp hoặc transaction hash. Unique constraint ngăn chart point và chi tiết giao dịch bị ghi nhiều lần. Foreign key với quy tắc cascade áp dụng cho dữ liệu thuộc người dùng như auth account, watchlist, alert và subscription. Các trường provider không bảo đảm luôn xuất hiện, như symbol, logo hoặc phần trăm biến động, được phép rỗng sau bước validation thay vì làm thất bại toàn bộ bản ghi.

Wallet address và token address có thể xuất hiện trước metadata tương ứng nên không bị buộc vào một bảng cha vật lý. Service chuẩn hóa định danh trước khi đọc hoặc ghi và sử dụng cùng khóa nghiệp vụ để ghép dữ liệu. Dữ liệu công khai từ blockchain được tách khỏi dữ liệu cá nhân; các truy vấn profile, watchlist, cảnh báo, chat và thanh toán luôn giới hạn theo user ID đã xác thực. Password được lưu dưới dạng hash, còn API key và secret của provider nằm trong cấu hình môi trường, không đi vào bảng nghiệp vụ.

% TODO(runtime): Hoàn thiện ghi alert_history sau khi delivery thành công.
% TODO(runtime): Bổ sung API lấy lịch sử, đánh dấu đã đọc và chống gửi trùng.

\subsection{Tổng kết thiết kế cơ sở dữ liệu}
\label{subsec:db_tong_ket}

Thiết kế cơ sở dữ liệu kết hợp quan hệ sở hữu chặt chẽ với các read model cache theo miền. Cách tổ chức này giữ toàn vẹn cho tài khoản, thanh toán và cảnh báo, đồng thời cho phép token và ví được đồng bộ theo address, thời gian và hạn mức provider. Bốn sơ đồ trong phần này trình bày đủ cấu trúc chính mà không đưa toàn bộ bảng cache vật lý vào một ERD quá rộng.
